\documentclass{beamer}
\usepackage[utf8]{inputenc}
\usetheme{Antibes}

% À chaque nouvelle section, on dit où on est
\AtBeginSection[ ]
{
 \begin{frame}<beamer>
   \frametitle{Plan}
   \tableofcontents[currentsection]
  \end{frame}
} 

\title{A Dynamic End-to-End Security for Coordinating Multiple Protections within a Linux Desktop}
\subtitle{COLSEC 2010}
\author{J.~Briffaut, M.~Peres \& C. Toinard}
\institute{ENSI de Bourges}

\begin{document}

	% page de garde
	\begin{frame}
		\titlepage
	\end{frame}

	% Ajoute le logo à toutes les autres pages
	\logo{\includegraphics[height=1.0cm]{figures/ensib_logo.png}}

	\section{Introduction}

		\begin{frame}
			\begin{block}{}
				Today's protection models are:
				\begin{itemize}
					\item mostly static
					\item made of several security components (firewalls, MAC, etc...)
					\item not adaptating to the user's task
				\end{itemize}
			\end{block}

			\begin{block}{}
				This leads administrators to:
				\begin{itemize}
					\item grant a user all the rights he needs, ever
					\item make use of virtual/physical machines dedicated for every single sensitive task
					\item think about hardware price and usability vs security
				\end{itemize}
			\end{block}
		\end{frame}

		\begin{frame}
			\begin{block}{}
				A protection model should:
				\begin{itemize}
					\item not require the use of virtualization
					\item be driven by the user's current task
					\item involve all the security components
				\end{itemize}
			\end{block}
		\end{frame}

		\begin{frame}
			\begin{block}{Example of security domains}
				\begin{itemize}
					\item email
					\item accounting/taxes
					\item epayment
				\end{itemize}
			\end{block}

			\begin{block}{Example of security components}
				\begin{itemize}
					\item firewalls (IPTables)
					\item loggers (Syslog)
					\item MAC systems (SELinux)
				\end{itemize}
			\end{block}
		\end{frame}

	\section{Architecture}
		\subsection*{Basics}
			\begin{frame}
				\begin{block}{}
					To implement such a protection model, it is needed to:
					\begin{itemize}
						\item determine which task the user in performing
						\item have a security model to match the user's task with
						\item glue all the security component and coordinate them
					\end{itemize}
				\end{block}
			\end{frame}

		\subsection*{Tracking the user}
			\begin{frame}
				\begin{block}{Tracking's requirement}
					Should gather information upon what the user is doing.\\
					Every action tracked should be overwritable.
				\end{block}

				\begin{block}{Determining what the user does}
					It requires:
					\begin{itemize}
						\item low-level information (sys-calls)
						\item high-level information (plugins)
					\end{itemize}
				\end{block}
			\end{frame}

		\subsection*{Having a security model}
			\begin{frame}
				\begin{block}{Security model's requirement}
					Should dynamically tell if the user is allowed to perform an action or not.\\
					This involves the creation of domains and transitions.
				\end{block}

				\begin{block}{}
					A security model is a finite-state machine where:
					\begin{itemize}
						\item State = Security Domain
						\item Transition = a user-generated event
					\end{itemize}
				\end{block}
			\end{frame}

		\subsection*{Gluing the security components}
			\begin{frame}
				\begin{block}{Glue's requirement}
					The security components should be drivable by the security model.\\
					They are to enforce the security model's decision
				\end{block}

				\begin{block}{Example of a browsing desktop}
					\begin{itemize}
						\item The user wants to browse http://www.google.com using Firefox.
						\item The finite state agrees with it.
						\item The firewall should open www.google.com:80.
						\item The MAC system should allow Firefox to connect to www.google.com:80.
					\end{itemize}
				\end{block}
			\end{frame}

	\section{Implementation}
		% Je donne des détails sur la structure, le daemon, les protocoles

	\section{Experimentation}
		% Je montre ici notre exemple, comme dans le papier et je le commente

	\section{Conclusion}
		% En quoi c'est mieux, en quoi ça respecte ce qu'on disait en intro

\end{document}